%!TEX TS-program = lualatex
%!TEX encoding = UTF-8 Unicode

\documentclass[11pt, addpoints]{exam}
\usepackage[left=1in,right=0.75in,top=1in]{geometry}     

\usepackage{fontspec}
\def\mainfont{Linux Libertine O}
\setmainfont[Ligatures=TeX, Contextuals={NoAlternate}, BoldFont={* Bold}, ItalicFont={* Italic}, BoldItalicFont={* BoldItalic}, Numbers={Proportional}]{\mainfont}
\setmonofont[Scale=MatchLowercase]{Inconsolatazi4} 
\setsansfont[Scale=MatchLowercase]{Linux Biolinum O} 
\usepackage{microtype}

\usepackage{graphicx}
	\graphicspath{%
	{/Users/goby/Pictures/teach/151/exams/}%
	{/img/}}%

% Used to get the semestesr and last two digits of the year for the header below.
\def\short#1{\csname @gobbletwo\expandafter\endcsname\number#1}
\def\Semester{\ifthenelse{\number\month>8}
	{F}%
	{%
		\ifthenelse{\number\month<4}
		{S}%
		{N}%
	}%
}

\pagestyle{headandfoot}
\firstpageheader{BI 151, Quizzam 3, \Semester{}\short{\year}}{}{%
	\ifprintanswers\textbf{KEY}\else Name: \enspace \makebox[2.5in]{\hrulefill}\fi}
\runningheader{}{}{\small{pg. \thepage}}

\footer{}{}{}
\extrafootheight{-0.5in}

\pointsinmargin
\marginpointname{ pts}

%\usepackage{multicol}
%\usepackage{booktabs}

%% Remove comment %% to print answer key.
%\printanswers
\unframedsolutions
\renewcommand{\solutiontitle}{}
\SolutionEmphasis{\bfseries}


\begin{document}

\fullwidth{\noindent Read each question completely. Explain \emph{thoroughly} your reasoning using clear and concise sentences. \vspace{1\baselineskip}}

\begin{questions}

\question[5]
Clearly explain what is meant by ``mosaic evolution'' or the ``mosaic
character of evolution'' and how it relates to transitional forms found
in the fossil record. Support your explanation with two \emph{specific}
traits observed in the reptile / bird fossils shown in class.

\ifprintanswers
	\vspace*{\stretch{1}}
\else
	\vspace*{\stretch{1}}
\fi

\question[10]
Describe two pieces of evidence \emph{from transitional forms} that
support the hypothesis that reptiles and mammals share a common
ancestor. Explain why each piece of evidence is considered to support
this hypothesis. Clearly explain how this evidence from the fossil
record is consistent with the geological time scale.

\ifprintanswers
	\vspace*{\stretch{1}}
\else
	\vspace*{\stretch{1}}
\fi

\newpage

\fullwidth{\noindent The following hypothesis makes specific \emph{predictions} about
transitional forms between organisms and about when the forms should
appear in the fossil record. \emph{Thoroughly explain your reasoning for
questions~\ref{ques:fossil_predictions} and \ref{ques:fossil_evidence}}. (MYA = millions of years ago)}
\begin{center}
	\includegraphics{exam3_phylo_tree}
\end{center}

\question[10]\label{ques:fossil_predictions}
Identify \emph{all} transitional forms that are \emph{predicted} by this hypothesis.
Based on this hypothesis, when is each transitional form
\emph{predicted} to appear in the fossil record?

\ifprintanswers
	\vspace*{\stretch{1}}
\else
	\vspace*{\stretch{1}}
\fi

\newpage

\question[10]\label{ques:fossil_evidence}
You find evidence in the fossil record that allows you to test the
above hypothesis. You find the following transitional forms in the
fossil record at the following geological ages: \vspace{\baselineskip}

fish/amphibian: 375 million years ago

reptile/mammal: 290 million years ago

first vertebrates: 490 million years ago

reptile/bird: 200 million years ago

amphibian/reptile: 325 million years ago\vspace{\baselineskip}

Thoroughly explain \emph{all} of the ways in which your new evidence
supports or falsifies the hypothesis shown on the previous page.

\ifprintanswers
	\vspace*{\stretch{1}}
\else
	\vspace*{\stretch{1}}
\fi

\newpage

\question[7]
Draw a new hypothesis that is consistent with \emph{all} of the new evidence
provided by question~\ref{ques:fossil_evidence}. Include a properly labeled time
scale.

\ifprintanswers
	\vspace*{\stretch{6}}
\else
	\vspace*{\stretch{6}}
\fi

\question[2]
Name the type of monomers that are joined together to form a
carbohydrate.

\ifprintanswers
	\textbf{Saccharides}\vspace*{\stretch{1}}
\else
	\vspace*{\stretch{1}}
\fi

\question[2]
Sucrose, a carbohydrate, is made of two monomers, glucose and
fructose. Name the type of reaction that occurs to \emph{break} the bond
that connects together the two monomers.

\ifprintanswers
	\textbf{hydrolysis}\vspace*{\stretch{1}}
\else
	\vspace*{\stretch{1}}
\fi

\question[2]
Does this reaction release or consume energy?

\ifprintanswers
	\textbf{Release}\vspace*{\stretch{1}}
\else
	\vspace*{\stretch{1}}
\fi

\question[2]
Name the group of macromolecules that are used for information storage. 

\ifprintanswers
	\textbf{Nucleic Acids}\vspace*{\stretch{1}}
\else
	\vspace*{\stretch{1}}
\fi

\end{questions}


\end{document}